% Volume II: Causal Explosion
% Part III. The Combinatorics of Context
% Chapter 13: Counterfactual Branching
% Spine: proof-spines/ch013-spine.md

\chapter{Counterfactual Branching}
\label{ch:ch013}

Responsibility depends on \textbf{counterfactual branching}: to say that an act mattered is to show that nearby interventions change the outcome, so that
\[
Y_{do(A=a)}\neq Y_{do(A=a')}.
\]
\Definition\ This intervention contrast fixes the book's difference-making criterion for saying that an act mattered.
Legal events, however, rarely involve a single isolated switch. They involve multiple agents, delayed consequences, feedback, strategic adaptation, and the familiar problems of structural causal models, potential outcomes, actual causation, overdetermination, preemption, and omissions.

To keep this branching disciplined, the chapter introduces the counterfactual branching functional
\[
\mathcal B(X)
=
\sum_{a\in\mathcal A}
|\Omega_{do(a)}|,
\]
\Definition\ This branching functional names the disciplined comparison class for nearby alternatives.
which counts the materially distinct outcome spaces generated by plausible interventions. The aim is not to indulge fantasy about unreal worlds, but to compare the actual path with nearby mechanism-grounded alternatives. Negligence, intent, prevention, and responsibility all depend on which branches are live, how many of them there are, and which changes would have made a difference.
